\SourceLeaf{302}
\par\bigskip\pdfbookmark[0]{Bemerkungen zu Über den Zusammenhang}{nexus-bemerkungen}
\begin{center}
\textsc{Bemerkungen zur Abhandlung}\\[1em]
\textsc{Über den Zusammenhang zwischen der Anzahl der Klassen, in die binäre quadratische Formen zerfallen, und ihrer Determinante.}
\end{center}
\BookRule{.16\textwidth}
\par\medskip\pdfbookmark[1]{Zu I. und II.}{nexus-note-I-II}\centerline{Zu I. und II.}\smallskip

Die zweite Formel für die Anzahl der Punkte innerhalb des Kreises (I, Artikel 3, und II, Artikel 5) ergibt sich aus der Betrachtung des in den Kreis eingeschriebenen Quadrats, dessen Seiten parallel zu den Koordinatenachsen sind; der Vergleich beider Formeln führt zu dem auch arithmetisch leicht zu beweisenden Satz
\[
r'+r''+\ldots+r^{(q)}=qq+r^{(q+1)}+r^{(q+2)}+\ldots+r^{(r)}
\]
aus dem sich wiederum die Richtigkeit der ersten der beiden folgenden Regeln ergibt, die sich auf einem besonderen Blatt vorfanden:

„Auflösungen der Gleichung $xx+yy\leqq A$; Formel
\[
1+4\surd A+4\surd\tfrac12 A+8\Sigma(\surd(A-nn)-n)
\]
wobei bei jeder Wurzel der Bruchteil weggelassen und von $n=1$ bis $n=\surd A$ „(soll heißen $\surd\tfrac12 A$)“ summiert wird.

Andere Formel
\[
1+4\{A-\tfrac A3+\tfrac A5-\tfrac A7+\tfrac A9-\tfrac A{11}\ldots\}
\]
wobei bei jedem Glied der Bruchteil weggelassen wird.“

Diese letztere Formel folgt aus dem später (I, Artikel 6) angewandten Satz über die Anzahl aller verschiedenen Darstellungen einer bestimmten Zahl durch die Form $xx+yy$ (vgl. Disqq. Arithm., Artikel 182, Anmerkung). Er lässt sich leicht folgendermaßen umformen: Die Anzahl der verschiedenen Darstellungen einer positiven ganzen Zahl $m$ durch die Form $xx+yy$ ist $=4(a-b)$, wobei $a,b$ die Anzahlen der Teiler von $m$ bezeichnen, die jeweils die Form $4n+1,\ 4n+3$ haben. Aus dem Vergleich
\SourceLeaf{303}
dieser arithmetischen Formel mit der (in I, Artikel 5, oder II, Artikel 4) durch geometrische Betrachtungen gewonnenen mittleren Darstellungsanzahl erhält man leicht und in aller Strenge das bekannte Ergebnis
\[
\frac\pi4=1-\tfrac13+\tfrac15-\tfrac17+\ldots
\]
das in der Abhandlung (I, Artikel 7) durch einen ähnlichen Vergleich, aber mithilfe unendlicher Produkte hergeleitet wird.
\BookRule{.25\textwidth}
